\documentclass[12pt]{article}
\newcommand{\ul}[1]{\underline{#1}}

\setlength{\textwidth}{6.5in}
\setlength{\textheight}{9.0in}
\setlength{\oddsidemargin}{-0.25in}
\setlength{\evensidemargin}{-0.25in}
\setlength{\topmargin}{-0.75in}

\begin{document}
\noindent{\large\bf PHYS-3350 \hfill Thermal Physics \hfill Spring 2026} \\
Homework 1 - Preliminaries \hfill {\bf Brandon Aikman}

% Homework 1 
%\begin{center}
%	{\bf Chapter n: ex}
%\end{center}
\begin{enumerate}
	% 1.
	\item What is the mass of 4 moles of carbon dioxide? \\
	\textbf {Given: } $m_c  = 12$ g, $m_o = 16$ g, \textbf{Find: } $m$.
	\begin{eqnarray*}
		m &=& 4 \cdot (m_c + (2 \cdot m_o))\\
		&=& 4 \cdot (12 \mbox{ g} + (2 \cdot 16 \mbox{ g}))\\
		&=& 4 \cdot 44 \mbox{ g}.
	\end{eqnarray*}
	$$\fbox{$= 176$ g}$$
	
	% 2.
    \item A typical bacterium has a mass of 1 pg. Calculate the mass of a mole of bacteria. \\
    \textbf {Given: } $m_b = 1$ pg, $m_e  = 5000$ kg, \textbf{Find: } $m$.
    \begin{eqnarray*}
		m &=& N_A \cdot m_b\\
		&=& 6.022 \cdot 10^{23} \cdot 1 \mbox{ pg}\\
		&=& 6.022 \cdot 10^{8} \mbox{ kg}.
	\end{eqnarray*}
	$$\fbox{$\approx 120000$ elephant-masses}$$
    
    % 3.
    \item Convert the following temperatures into Kelvin:
    \begin{enumerate}
        \item Boiling point of hydrogen: -252.87 degrees Celsius.
        \begin{eqnarray*}
			-252.87\mbox{ $^{ \circ}$C} &=& 273.15 + (-252.87) \mbox{ K}.
		\end{eqnarray*}
		$$\fbox{$= 20.28$ K}$$
        \item Triple point of water: 0.01 degrees Celsius.
        \begin{eqnarray*}
			0.01\mbox{ $^{ \circ}$C} &=& 273.15 + (0.01) \mbox{ K}.
		\end{eqnarray*}
		$$\fbox{$= 273.16$ K}$$
    \end{enumerate}
    
    % 4.
    \item Convert the following temperatures into Fahrenheit:
    \begin{enumerate}
        \item -13 degrees Celsius.
        \begin{eqnarray*}
			-13\mbox{ $^{ \circ}$C} &=& -13 \cdot 1.8 + 32 \mbox{ $^{ \circ}$F}.
		\end{eqnarray*}
		$$\fbox{$= 8.6 \mbox{ $^{ \circ}$F}$}$$
        \item 25 degrees Kelvin.
        \begin{eqnarray*}
        	25 \mbox{ K} &=& 25 - 273.15 \mbox{ $^{ \circ}$C}\\
        	&=& -248.15 \mbox{ $^{ \circ}$C}\\
        	-248.15 \mbox{ $^{ \circ}$C} &=& -248.15 \cdot 1.8 + 32 \mbox{ $^{ \circ}$F}\\
			&=& -414.67 \mbox{ $^{ \circ}$F}
		\end{eqnarray*}
		$$\fbox{$\approx -410 \mbox{ $^{ \circ}$F}$}$$
    \end{enumerate}
    
    % 5.
    \item A system contains $n$ atoms, each of which can only have zero or one quanta of energy. How many ways can you arrange $r$ quanta of energy when:
    \begin{enumerate}
        \item $n = 2$, $r = 1$
        	$${2 \choose 1} = \fbox{2}.$$
        \item $n = 20$, $r = 10$
        	$${20 \choose 10} = \fbox{184756}.$$
        \item $n = 2 \cdot 10^{23}$, $r = 10^{23}$
        	$$\fbox{${2 \cdot 10^{23} \choose 10^{23}}$}.$$
    \end{enumerate}
\end{enumerate}

\end{document}
